\chapter{Dalla frequenza alla probabilità}

Nel capitolo precedente abbiamo visto che, quando i risultati di un esperimento sono «alla pari», la frequenza
relativa di un evento si avvicina sempre di più al numero $\card{A}/\card{\Omega}$. È il momento di dare un nome a
questo numero: la \textbf{probabilità}. In questo capitolo la definiamo come il valore a cui tende la frequenza,
la calcoliamo in alcuni casi interessanti (le mani del poker, un dado truccato) e ne ricaviamo le regole di calcolo:
la probabilità del complementare, dell'unione, della differenza, la \textbf{probabilità condizionata}, la legge della
\textbf{probabilità totale} e l'\textbf{indipendenza}.

% ============================================================
\section{La probabilità come limite della frequenza}
% ============================================================

La frase «la frequenza si avvicina sempre di più a un numero» ha un nome preciso in matematica: \emph{limite}. Prima
di usarla, vediamo che cosa significa esattamente.

\begin{richiamo}{il limite di una successione}
Una \term{successione} è un elenco infinito di numeri, uno per ogni naturale $n$: $a_1, a_2, a_3, \ldots$ Per
esempio $a_n = 1/n$ è la successione $1, \tfrac12, \tfrac13, \tfrac14, \ldots$, e la frequenza $f_n(A)$ di un evento,
calcolata dopo $n$ prove, è una successione.

Si dice che la successione \term{tende} al numero $L$ (o che ha \term{limite} $L$), e si scrive
\[ \lim_{n \to \infty} a_n = L \qquad\text{oppure}\qquad a_n \to L \ \text{per}\ n \to \infty , \]
se i suoi termini si avvicinano a $L$ \textbf{quanto si vuole} e \textbf{per sempre}. In modo preciso: per quanto
stretta si scelga una fascia attorno a $L$, cioè per qualsiasi tolleranza $\varepsilon > 0$, da un certo indice $N$
in poi \emph{tutti} i termini stanno nella fascia:
\[ \abs{a_n - L} < \varepsilon \qquad \text{per ogni } n \geq N . \]
Il simbolo $\abs{a_n - L}$ è la distanza fra $a_n$ e $L$ (il valore assoluto della differenza). Più stretta è la
fascia, più avanti bisogna andare per trovare $N$, ma un $N$ esiste sempre.
\begin{center}
\begin{tikzpicture}
\begin{axis}[quaderno, width=11cm, height=5.2cm, xmin=0, xmax=31, ymin=-0.7, ymax=2.7,
  xtick={1,5,8,15,20,25,30}, ytick={0,1,2}, yticklabels={$0$,$L = 1$,$2$},
  xlabel=$n$, ylabel=$a_n$, grid=none]
  \fill[evid, opacity=0.6] (axis cs:0,0.8) rectangle (axis cs:31,1.2);
  \node[font=\scriptsize, anchor=south east] at (axis cs:31,1.2) {fascia da $L - \varepsilon$ a $L + \varepsilon$};
  \draw[mspiB, dashed, line width=0.8pt] (axis cs:0,1) -- (axis cs:31,1);
  \draw[mspigray, dashed] (axis cs:8,-0.7) -- (axis cs:8,2.7) node[above, font=\scriptsize] {$N = 8$};
  \addplot[only marks, mark=*, mark size=1.7pt, mspiA, samples at={1,2,...,30}] {1 + ifthenelse(mod(x,2)==0,1,-1)*1.5/x};
\end{axis}
\end{tikzpicture}
\end{center}
Nel disegno $a_n = 1 + (-1)^n \cdot 1{,}5/n$: i termini saltano sopra e sotto il valore $L = 1$, ma sempre più vicino.
Con la tolleranza $\varepsilon = 0{,}2$ (la fascia gialla) basta prendere $N = 8$: da lì in poi nessun termine esce
più. Con una fascia più stretta servirebbe un $N$ più grande, ma lo si troverebbe comunque.

Altri esempi: $1/n \to 0$; $(n+1)/n = 1 + 1/n \to 1$. Invece la successione $(-1)^n = -1, 1, -1, 1, \ldots$
\textbf{non} ha limite: salta per sempre fra due valori e non si avvicina a nessun numero.
\end{richiamo}

Ora possiamo dare la definizione che regge tutta la prima parte del corso.

\dfn{Probabilità (definizione frequentista)}{
  Dato uno spazio campione discreto $\Omega$ e un suo sottoinsieme (evento) $A$, la \term{probabilità} che si
  verifichi $A$ è il limite della frequenza relativa di $A$ quando il numero di prove tende all'infinito:
  \[ P(A) = \lim_{n \to \infty} \frac{n_A}{n} . \]
}

\begin{figure}[H]
\centering
\begin{tikzpicture}
\begin{axis}[width=13cm, height=6cm, xmode=log, xmin=1, xmax=5000, ymin=0, ymax=1,
  xlabel={numero di lanci $n$ (scala logaritmica)}, ylabel={$f_n(\text{Testa})$},
  log ticks with fixed point, xtick={1,10,100,1000,5000}, ytick={0,0.25,0.5,0.75,1},
  yticklabels={$0$,$0{,}25$,$0{,}5$,$0{,}75$,$1$}, tick label style={font=\scriptsize}, label style={font=\small},
  grid=major, grid style={mspigray!20}, axis lines=left, legend pos=north east, legend style={font=\scriptsize}]
  \draw[mspigray, dashed, line width=0.9pt] (axis cs:1,0.5) -- (axis cs:5000,0.5);
  \addplot[mspiA, line width=0.9pt] coordinates {(1,1.0000) (2,0.5000) (3,0.6667) (4,0.5000) (5,0.4000) (6,0.5000) (7,0.5714) (8,0.5000) (9,0.5556) (10,0.6000) (12,0.5833) (14,0.5714) (16,0.5625) (18,0.5000) (20,0.4500) (23,0.4348) (26,0.4615) (30,0.4333) (35,0.4286) (40,0.4500) (45,0.4444) (50,0.4400) (60,0.3833) (70,0.3857) (80,0.4125) (90,0.4222) (100,0.4300) (115,0.4174) (130,0.4154) (150,0.4533) (175,0.4686) (200,0.4800) (230,0.4696) (260,0.4846) (300,0.4833) (350,0.4829) (400,0.4900) (450,0.4889) (500,0.4900) (600,0.4900) (700,0.5057) (800,0.5050) (900,0.5067) (1000,0.5090) (1200,0.5108) (1400,0.5086) (1600,0.5169) (1800,0.5133) (2000,0.5060) (2500,0.5052) (3000,0.4997) (3500,0.5049) (4000,0.5038) (5000,0.5098)};
  \addplot[mspiB, line width=0.9pt] coordinates {(1,1.0000) (2,0.5000) (3,0.3333) (4,0.5000) (5,0.4000) (6,0.3333) (7,0.4286) (8,0.3750) (9,0.3333) (10,0.3000) (12,0.4167) (14,0.4286) (16,0.4375) (18,0.3889) (20,0.3500) (23,0.3913) (26,0.4615) (30,0.5000) (35,0.4857) (40,0.4500) (45,0.4889) (50,0.4800) (60,0.5000) (70,0.5286) (80,0.5500) (90,0.5333) (100,0.5500) (115,0.5478) (130,0.5385) (150,0.5467) (175,0.5143) (200,0.4950) (230,0.4913) (260,0.5154) (300,0.4967) (350,0.4829) (400,0.4900) (450,0.4844) (500,0.4780) (600,0.4733) (700,0.4800) (800,0.4775) (900,0.4844) (1000,0.4800) (1200,0.4825) (1400,0.5036) (1600,0.5081) (1800,0.5050) (2000,0.5080) (2500,0.5032) (3000,0.5043) (3500,0.5086) (4000,0.5038) (5000,0.5044)};
  \addplot[mspiC, line width=0.9pt] coordinates {(1,0.0000) (2,0.5000) (3,0.6667) (4,0.7500) (5,0.6000) (6,0.5000) (7,0.4286) (8,0.5000) (9,0.5556) (10,0.6000) (12,0.5833) (14,0.6429) (16,0.5625) (18,0.5556) (20,0.5500) (23,0.5217) (26,0.5385) (30,0.5333) (35,0.4857) (40,0.4250) (45,0.4667) (50,0.5000) (60,0.5167) (70,0.5000) (80,0.5125) (90,0.5111) (100,0.5000) (115,0.4957) (130,0.4769) (150,0.4933) (175,0.4971) (200,0.5150) (230,0.5043) (260,0.5077) (300,0.5033) (350,0.4914) (400,0.4925) (450,0.4889) (500,0.4960) (600,0.4933) (700,0.4971) (800,0.4888) (900,0.4867) (1000,0.4800) (1200,0.4867) (1400,0.4900) (1600,0.4869) (1800,0.4856) (2000,0.4845) (2500,0.4884) (3000,0.4903) (3500,0.4923) (4000,0.4900) (5000,0.4944)};
  \legend{1ª serie, 2ª serie, 3ª serie}
\end{axis}
\end{tikzpicture}
\caption{Frequenza di «testa» in tre serie di 5000 lanci di una moneta onesta (simulate). All'inizio le tre curve
oscillano moltissimo e sono diverse fra loro; al crescere di $n$ si stringono tutte attorno a $1/2$.}
\end{figure}

\warning{Un limite un po' speciale}{
  La frequenza non è una successione come $1/n$, scritta una volta per tutte: ogni serie di lanci produce numeri
  diversi (nella figura le tre curve non coincidono mai). Il limite va quindi inteso in senso «fisico»: è il valore
  attorno a cui \textbf{tutte} le serie finiscono per stringersi. Rendere preciso questo discorso è il compito della
  \emph{legge dei grandi numeri}, che vedremo nella terza parte del corso; il capitolo seguente aggira il problema con
  l'approccio assiomatico. Per ora la definizione frequentista basta per capire da dove vengono tutte le regole.
}

\subsection{Spazi finiti con eventi equiprobabili}

Se gli eventi elementari sono equivalenti, nel senso del capitolo precedente, ora possiamo dire che sono
\term{equiprobabili}: hanno tutti la stessa probabilità $1/\card{\Omega}$. Per uno spazio finito con eventi
elementari equiprobabili vale allora la relazione generale
\[ P(A) = \frac{\card{\Omega_A}}{\card{\Omega}} , \]
dove $\Omega_A$ è l'insieme di tutti gli esiti che comportano il verificarsi di $A$ (cioè l'insieme $A$ stesso).

\begin{formula}[spazi equiprobabili]
\[ P(A) = \frac{\card{A}}{\card{\Omega}} = \frac{\text{numero dei casi favorevoli}}{\text{numero dei casi possibili}} \]
\end{formula}

Con questa formula si rileggono tutti gli esempi già visti: la probabilità che un dado onesto dia un numero pari è
$3/6 = 1/2$; quella di ottenere almeno un 6 in due lanci è $11/36$; quella di indovinare un PIN a caso è $1/10\,000$.
Il punto delicato è sempre lo stesso: i risultati che si contano devono essere \textbf{davvero} equiprobabili.

\begin{esempio}{L'errore di d'Alembert}
Si lanciano due monete. Qual è la probabilità di ottenere una testa e una croce?

\emph{Ragionamento sbagliato}: i risultati sono tre, «due teste», «due croci», «una testa e una croce», quindi la
probabilità è $1/3$. Lo sostenne nel Settecento il matematico Jean-Baptiste d'Alembert.

\emph{Ragionamento giusto}: i risultati equiprobabili sono le quattro coppie ordinate $TT$, $TC$, $CT$, $CC$ (lo si
vede pensando alle due monete come diverse, Capitolo~1). L'evento «una testa e una croce» è $\{TC, CT\}$ e contiene
\textbf{due} di esse:
\[ P(\{TC, CT\}) = \frac{2}{4} = \frac{1}{2} . \]
I tre risultati di d'Alembert esistono, ma non sono equiprobabili: «una testa e una croce» si può ottenere in due
modi, gli altri in uno solo. Basta lanciare due monete qualche centinaio di volte per vedere che la frequenza si
avvicina a $1/2$, non a $1/3$.
\end{esempio}

% ============================================================
\section{Un esempio completo: i punteggi del poker}
% ============================================================

Il poker «all'italiana» si gioca con un mazzo ridotto: le carte dal 7 all'asso, cioè 8 valori (7, 8, 9, 10, J, Q, K,
A) per 4 semi, per un totale di $8 \times 4 = 32$ carte. Un giocatore riceve 5 carte, estratte senza reinserimento.

\begin{figure}[H]
\centering
\begin{tikzpicture}[x=0.95cm, y=0.82cm, carta/.style={draw=mspigray, rounded corners=1.5pt, minimum width=0.78cm,
                    minimum height=0.66cm, inner sep=0pt, font=\small}]
  \foreach \v [count=\c] in {7,8,9,10,J,Q,K,A}
    \node[font=\scriptsize\bfseries, mspigray] at (\c,0.8) {\v};
  \foreach \s/\col [count=\r] in {\heartsuit/red!75!black, \diamondsuit/red!75!black, \clubsuit/black, \spadesuit/black} {
    \foreach \v [count=\c] in {7,8,9,10,J,Q,K,A}
      \node[carta, text=\col] (c\r\c) at (\c,-\r+1) {\v$\s$};
  }
  % una mano: coppia d'assi
  \foreach \r/\c in {1/8, 4/8, 2/1, 3/3, 1/7}
    \node[carta, fill=evid, draw=evid!50!black, line width=0.8pt] at (\c,-\r+1) {};
  \foreach \s/\col [count=\r] in {\heartsuit/red!75!black, \diamondsuit/red!75!black, \clubsuit/black, \spadesuit/black}
    \foreach \v [count=\c] in {7,8,9,10,J,Q,K,A}
      \node[text=\col, font=\small] at (\c,-\r+1) {\v$\s$};
  \node[font=\small, anchor=west, text width=4.3cm] at (8.7,-1.5)
    {In giallo una mano con una \textbf{coppia d'assi}: $A\heartsuit$, $A\spadesuit$, $7\diamondsuit$, $9\clubsuit$,
     $K\heartsuit$.};
\end{tikzpicture}
\caption{Il mazzo da 32 carte: 8 valori per 4 semi (cuori, quadri, fiori, picche).}
\end{figure}

In tutti i calcoli lo spazio campione è l'insieme delle \textbf{cinquine} di carte distinte, in cui l'ordine non
conta: sono i sottoinsiemi di 5 carte di un insieme di 32. Se il mazzo è ben mescolato, sono tutti equiprobabili, e
\[ \card{\Omega} = \binom{32}{5} = \frac{32!}{5!\,27!} = \frac{32 \cdot 31 \cdot 30 \cdot 29 \cdot 28}{5 \cdot 4 \cdot
   3 \cdot 2} = 201\,376 . \]
Ogni probabilità sarà quindi un rapporto con denominatore $201\,376$: tutto il lavoro sta nel contare il numeratore.

\subsection{Coppia d'assi e coppia qualsiasi}

\begin{esempio}{Due assi (anche con punteggi più alti)}
Sia $C_2$ l'evento «nella mano ci sono esattamente due assi». Una cinquina di $C_2$ si costruisce in due passi:
si mettono i 4 assi del mazzo su 2 posti ($\binom{4}{2}$ modi) e si sceglie una terna fra le $32 - 4 = 28$ carte che
non sono assi ($\binom{28}{3}$ modi). Per il principio fondamentale:
\[
  \card{C_2} = \binom{4}{2}\binom{28}{3} = \frac{4!}{2!\,2!} \cdot \frac{28!}{3!\,25!} = 6 \cdot 3\,276 = 19\,656 ,
  \qquad
  P(C_2) = \frac{\card{C_2}}{\card{\Omega}} = \frac{19\,656}{201\,376} \approx 0{,}0976 .
\]
Così però si contano anche mani migliori di una semplice coppia d'assi: il \emph{full} (due assi e tre carte uguali
fra loro) e la \emph{doppia coppia} (due assi e un'altra coppia), perché nella terna le carte possono avere lo stesso
valore.
\end{esempio}

\begin{esempio}{Solo una coppia d'assi}
Per escludere i punteggi più alti, le tre carte che non sono assi devono avere tre valori \textbf{diversi}. Restano
ferme le posizioni dei due assi ($\binom{4}{2}$ modi); poi:
\begin{itemize}
  \item la terza carta può essere una qualsiasi delle 28 che non sono assi;
  \item la quarta può essere solo una delle 24 che hanno un valore diverso dalla terza: delle 27 carte non-asso
        rimaste si tolgono le 3 dello stesso valore della terza;
  \item la quinta, per lo stesso motivo, solo una delle 20 con valore diverso dalle due precedenti.
\end{itemize}
In questo modo le tre carte vengono scelte \emph{in ordine}, mentre nella mano l'ordine non conta: ogni terna è stata
contata $3! = 6$ volte, e bisogna dividere.
\[
  \card{C_2'} = \binom{4}{2} \cdot \frac{28 \cdot 24 \cdot 20}{3!} = \frac{4 \cdot 3}{2} \cdot \frac{28 \cdot 24 \cdot
  20}{6} = 28 \cdot 24 \cdot 20 = 13\,440 ,
  \qquad
  P(C_2') = \frac{13\,440}{201\,376} \approx 0{,}067 .
\]
\textbf{Verifica con un altro conteggio.} Le tre carte non-asso si possono scegliere anche così: prima i tre valori
diversi fra i 7 valori che non sono l'asso ($\binom{7}{3} = 35$ modi), poi il seme di ciascuna ($4^3 = 64$ modi). Si
ottiene $35 \cdot 64 = 2\,240$ terne, e $6 \cdot 2\,240 = 13\,440$: lo stesso numero.
\end{esempio}

Se la coppia può essere di qualsiasi valore, la probabilità precedente va moltiplicata per 8: le coppie di 7 (quattro
sette su due posti), di 8, \ldots, di re si contano esattamente come quella d'assi, e i valori sono otto. La
probabilità di avere una coppia (e niente di meglio) è circa $8 \cdot 0{,}067 \approx 0{,}53$: più di una mano su due.

\subsection{Almeno tre assi, tris, colore}

\begin{esempio}{Almeno tre assi}
L'evento $C_3$ = «almeno tre assi» è l'unione di due eventi incompatibili: «esattamente tre assi» (i 4 assi su 3
posti, le altre 2 carte fra le 28 non-asso) ed «esattamente quattro assi» (i 4 assi su 4 posti, l'ultima carta fra le
28). Siccome non hanno mani in comune, si sommano:
\[
  \card{C_3} = \binom{4}{3}\binom{28}{2} + \binom{4}{4}\binom{28}{1} = 4 \cdot 378 + 1 \cdot 28 = 1\,540 ,
  \qquad
  P(C_3) = \frac{1\,540}{201\,376} \approx 0{,}0076 .
\]
\end{esempio}

\begin{esempio}{Un tris qualsiasi}
Sia $C_3'$ l'evento «un tris, con le altre due carte diverse» (quindi non un full e non un poker). Indichiamo con
$C_3'(7)$ le mani con esattamente un tris di 7. I tris di 8, 9, \ldots, assi si contano allo stesso modo, quindi
$\card{C_3'} = 8\,\card{C_3'(7)}$.

Per $C_3'(7)$: tre posizioni sono occupate da un 7, e i tre 7 si scelgono fra i quattro in $\binom{4}{3} =
\frac{4!}{3!\,1!} = 4$ modi; la quarta carta si sceglie fra le 28 che non sono 7, la quinta fra le 24 di valore
diverso dai 7 e dalla quarta. Le ultime due sono state scelte in ordine, quindi si divide per $2! = 2$:
\[
  \card{C_3'(7)} = 4 \cdot \frac{28 \cdot 24}{2} = 1\,344 ,
  \qquad
  P(C_3') = 8 \cdot \frac{1\,344}{201\,376} = 8 \cdot 0{,}006674 \approx 0{,}0534 .
\]
\end{esempio}

\begin{esempio}{Colore}
Sia $C_p$ l'evento «colore di picche», cioè cinque carte di picche. Ci sono 8 carte di picche da distribuire su 5
posti, con ordine inessenziale:
\[
  \card{C_p} = \binom{8}{5} = \frac{8 \cdot 7 \cdot 6}{3!} = 56 ,
  \qquad
  P(C_p) = \frac{56}{201\,376} \approx 0{,}00028 .
\]
(Si è usata la simmetria $\binom{8}{5} = \binom{8}{3}$.) La probabilità di un colore qualsiasi è $4 \cdot P(C_p)
\approx 0{,}0011$. Questo conteggio include però le \textbf{scale reali}, cioè le cinque carte consecutive dello stesso
seme, che valgono di più. Per ogni seme ce ne sono 5, dalla minima ($A, 7, 8, 9, 10$, con l'asso che vale come carta
più bassa) alla massima ($10, J, Q, K, A$). Togliendole:
\[
  \card{C_p} = 56 - 5 = 51 ,
  \qquad
  P(C_p) = \frac{51}{201\,376} \approx 0{,}00025 ,
\]
e la probabilità di un colore qualsiasi senza scale reali è $4 \cdot 0{,}00025 \approx 0{,}001$.
\end{esempio}

\begin{table}[H]
\centering
\renewcommand{\arraystretch}{1.25}
\begin{tabular}{@{}lrr@{}}
  \toprule
  \textbf{Mano} & \textbf{Casi favorevoli} & \textbf{Probabilità} \\
  \midrule
  una coppia (e niente di meglio) & $8 \cdot 13\,440 = 107\,520$ & $0{,}534$ \\
  tris (non full né poker) & $8 \cdot 1\,344 = 10\,752$ & $0{,}053$ \\
  almeno tre assi & $1\,540$ & $0{,}0076$ \\
  colore (senza scale reali) & $4 \cdot 51 = 204$ & $0{,}001$ \\
  \bottomrule
\end{tabular}
\caption{Alcune mani del poker con 32 carte, su $\card{\Omega} = 201\,376$ mani possibili. Nel poker vale di più la
mano più rara: un colore è circa cinquanta volte più raro di un tris.}
\end{table}

% ============================================================
\section{Quando i risultati non sono equiprobabili}
% ============================================================

Togliamo ora l'ipotesi che gli eventi elementari siano equiprobabili. Sia $\Omega$ uno spazio campione discreto
(finito o numerabile). La definizione di probabilità data prima vale ancora:
\begin{enumerate}[label=\alph*)]
  \item $A \subseteq \Omega$ è un evento e si conducono $n$ esperimenti;
  \item $n_A = n_A(n)$ è il numero di volte in cui $A$ si verifica;
  \item frequenza relativa e probabilità sono
        \[ f_n(A) = \frac{n_A}{n} , \qquad P(A) = \lim_{n \to \infty} f_n(A) = \lim_{n \to \infty} \frac{n_A}{n} . \]
\end{enumerate}
La differenza con il caso precedente è che ora \textbf{non} è più vero, in generale, che
$n_A \simeq n \card{A} / \card{\Omega}$. Quando gli eventi elementari non sono equiprobabili, un evento ha una
«misura» diversa da quella ordinaria, data dal numero dei suoi elementi: due eventi con lo stesso numero di elementi
($\card{A} = \card{B}$) possono avere pesi diversi, e le tecniche di conteggio non bastano più a calcolare le
probabilità.

\begin{esempio}{Un dado truccato}
Un dado è truccato in modo che il risultato 1 esca \textbf{5 volte su 6}; le altre cinque facce sono ugualmente
probabili fra loro. Si lancia una volta, quindi $\Omega = \{1, 2, 3, 4, 5, 6\}$. Siano $A = \{2, 4, 6\}$ (pari) e
$B = \{1, 3, 5\}$ (dispari): hanno entrambi 3 elementi.

Se $n_i$ è il numero di volte (su $n$ prove) in cui esce la faccia $i$, il trucco dice che l'1 esce cinque volte più
spesso di tutte le altre facce messe insieme:
\[ \sum_{i=1}^{6} n_i = n , \qquad n_1 \simeq 5 \sum_{i=2}^{6} n_i
   \quad\Longrightarrow\quad f_n(1) = \frac{n_1}{n} \simeq 5 \sum_{i=2}^{6} f_n(i) . \]
Le frequenze di tutte le facce sommano a 1 (dividendo la prima uguaglianza per $n$), quindi
\[
  f_n(1) + \sum_{i=2}^{6} f_n(i) = 1
  \quad\Longrightarrow\quad
  f_n(1) + \frac{1}{5} f_n(1) = f_n(1) \left( 1 + \frac{1}{5} \right) = 1
  \quad\Longrightarrow\quad
  f_n(1) \simeq \frac{5}{6} .
\]
Il restante $1 - \frac56 = \frac16$ si divide in parti uguali fra le altre cinque facce:
$f_n(2) = f_n(3) = \cdots = f_n(6) \simeq \dfrac{1 - 5/6}{5} = \dfrac{1}{30}$. Passando al limite:
\[ P(1) = \frac{5}{6} , \qquad P(i) = \frac{1}{30} \quad i = 2, 3, 4, 5, 6 . \]
Gli eventi $A$ e $B$ si verificano quando esce una delle loro facce, quindi le frequenze si sommano:
\[
  f_n(A) = \frac{n_2 + n_4 + n_6}{n} = f_n(2) + f_n(4) + f_n(6) \;\longrightarrow\; \frac{3}{30} = \frac{1}{10} ,
\]
\[
  f_n(B) = \frac{n_1 + n_3 + n_5}{n} = f_n(1) + f_n(3) + f_n(5) \;\longrightarrow\; \frac{5}{6} + \frac{2}{30} =
  \frac{9}{10} .
\]
In questo caso $P(B)$ ha una misura molto maggiore di $P(A)$, nove volte superiore, anche se i due eventi hanno lo
stesso numero di elementi.
\end{esempio}

\begin{figure}[H]
\centering
\begin{tikzpicture}
\begin{groupplot}[group style={group size=2 by 1, horizontal sep=1.6cm},
  width=6.4cm, height=4.6cm, ymin=0, ymax=0.95, xmin=0.3, xmax=6.7, xtick={1,...,6},
  ytick={0,0.0333,0.1667,0.8333}, yticklabels={$0$,$\frac{1}{30}$,$\frac16$,$\frac56$},
  tick label style={font=\scriptsize}, title style={font=\small}, axis lines=left,
  grid=major, grid style={mspigray!15}, xmajorgrids=false]
\nextgroupplot[title={dado onesto}]
  \addplot[bastoncini, mspiA] coordinates {(1,0.1667) (2,0.1667) (3,0.1667) (4,0.1667) (5,0.1667) (6,0.1667)};
\nextgroupplot[title={dado truccato}]
  \addplot[bastoncini, mspiB] coordinates {(1,0.8333) (2,0.0333) (3,0.0333) (4,0.0333) (5,0.0333) (6,0.0333)};
\end{groupplot}
\end{tikzpicture}
\caption{Le probabilità delle sei facce. Nel dado onesto si distribuiscono in parti uguali; in quello truccato quasi
tutta la probabilità sta sull'1.}
\end{figure}

% ============================================================
\section{Le proprietà della probabilità}
% ============================================================

Le frequenze relative obbediscono a regole semplicissime, che si ricavano contando. Passando al limite, le stesse
regole valgono per le probabilità. Per «passare al limite» serve sapere come si comportano i limiti con le operazioni.

\begin{richiamo}{le regole dei limiti}
Se due successioni hanno limite, $a_n \to a$ e $b_n \to b$, allora:
\[
  a_n + b_n \to a + b , \qquad a_n - b_n \to a - b , \qquad c \cdot a_n \to c \cdot a \ \ (c \text{ costante}) ,
  \qquad a_n \cdot b_n \to a \cdot b ,
\]
e, se $b \neq 0$, anche $a_n / b_n \to a / b$. In parole: il limite di una somma è la somma dei limiti, e così via.
Inoltre una successione costante ($a_n = c$ per ogni $n$) tende a $c$. Infine, se due successioni sono sempre uguali,
hanno lo stesso limite: un'uguaglianza fra frequenze, vera per ogni $n$, diventa un'uguaglianza fra probabilità.
\end{richiamo}

Siano $A$ e $B$ due eventi, $n_A$ e $n_B$ i loro numeri di occorrenze su $n$ prove e $f_n(A)$, $f_n(B)$ le frequenze
relative (entrambe $\geq 0$, chiaramente).

\teorema[proprieta]{Proprietà della frequenza e della probabilità}{
  Per due eventi qualsiasi $A$ e $B$:
  \begin{enumerate}[label=\alph*)]
    \item \textbf{eventi complementari}: $P(\compl{A}) = 1 - P(A)$;
    \item \textbf{unione}: $P(A \cup B) = P(A) + P(B) - P(A \cap B)$;
    \item \textbf{differenza}: $P(A \setminus B) = P(A) - P(A \cap B)$;
    \item \textbf{evento certo e impossibile}: $P(\Omega) = 1$ e $P(\emptyset) = 0$.
  \end{enumerate}
  Inoltre $0 \leq P(A) \leq 1$ per ogni evento $A$.
}

\begin{dimostrazione}
\passo{Complementare} In ogni prova o si verifica $A$ o si verifica $\compl{A}$, mai entrambi: le prove in cui si
verifica $\compl{A}$ sono tutte le altre, $n - n_A$. Quindi
\[ f_n(\compl{A}) = \frac{n - n_A}{n} = 1 - f_n(A) \quad\Longrightarrow\quad P(\compl{A}) = 1 - P(A) . \]
\passo{Unione} Se si sommano $n_A$ e $n_B$, le prove in cui si sono verificati \emph{entrambi} gli eventi (quelle di
$A \cap B$) vengono contate due volte, una in $n_A$ e una in $n_B$. Per contarle una volta sola bisogna toglierle una
volta:
\[
  f_n(A \cup B) = \frac{n_{A \cup B}}{n} = \frac{n_A + n_B - n_{A \cap B}}{n} = f_n(A) + f_n(B) - f_n(A \cap B)
\]
e, passando al limite, $P(A \cup B) = P(A) + P(B) - P(A \cap B)$.
\passo{Differenza} Perché si verifichi $A$ ma non $B$, dalle $n_A$ prove in cui si verifica $A$ vanno tolte quelle
in cui si verificano entrambi:
\[
  f_n(A \setminus B) = f_n(A \cap \compl{B}) = \frac{n_A - n_{A \cap B}}{n} = f_n(A) - f_n(A \cap B)
  \quad\Longrightarrow\quad P(A \setminus B) = P(A) - P(A \cap B) .
\]
\passo{Certo e impossibile} $\Omega$ si verifica in tutte le prove, $\emptyset$ in nessuna:
\[ f_n(\Omega) = \frac{n}{n} = 1 \ \Longrightarrow\ P(\Omega) = 1 , \qquad
   f_n(\emptyset) = \frac{0}{n} = 0 \ \Longrightarrow\ P(\emptyset) = 0 . \]
\passo{Fra 0 e 1} Per ogni $n$ vale $0 \leq n_A \leq n$, quindi $0 \leq f_n(A) \leq 1$, e il limite di numeri compresi
fra 0 e 1 è ancora compreso fra 0 e 1.
\end{dimostrazione}

\begin{figure}[H]
\centering
\begin{tikzpicture}
  \draw[mspigray, line width=0.5pt, rounded corners=2pt] (-2.7,-1.55) rectangle (2.7,1.55);
  \fill[mspiA!14] (-0.75,0) circle (1.2);
  \fill[mspiB!14] (0.75,0) circle (1.2);
  \begin{scope}\clip (-0.75,0) circle (1.2); \fill[evid] (0.75,0) circle (1.2); \end{scope}
  \draw[mspiA!85!black, line width=0.8pt] (-0.75,0) circle (1.2);
  \draw[mspiB!85!black, line width=0.8pt] (0.75,0) circle (1.2);
  \node[mspiA!85!black] at (-2.05,1.2) {$A$}; \node[mspiB!85!black] at (2.05,1.2) {$B$};
  \node[font=\small, align=center] at (-1.25,0) {solo $A$\\ 30};
  \node[font=\small, align=center] at (0,0) {$A \cap B$\\ 10};
  \node[font=\small, align=center] at (1.25,0) {solo $B$\\ 20};
  \node[font=\scriptsize, mspigray, anchor=south east] at (2.7,-1.55) {né $A$ né $B$: 40};
  \node[font=\small, anchor=west, text width=6.6cm] at (3.0,0)
    {Su $n = 100$ prove: $A$ si è verificato $n_A = 30 + 10 = 40$ volte, $B$ si è verificato $n_B = 10 + 20 = 30$
     volte. Ma $A \cup B$ si è verificato $30 + 10 + 20 = 60$ volte, non $40 + 30 = 70$: le 10 prove di $A \cap B$
     sono state contate due volte. Infatti $60 = 40 + 30 - 10$.};
\end{tikzpicture}
\caption{Perché nell'unione si sottrae l'intersezione. In ogni regione è scritto quante volte si è verificata.}
\end{figure}

La proprietà b) è nota anche come \term{sub-additività}: siccome $P(A \cap B) \geq 0$, ne segue
$P(A \cup B) \leq P(A) + P(B)$, con il segno di uguale esattamente quando $P(A \cap B) = 0$, per esempio se $A$ e
$B$ sono \textbf{incompatibili}. Per gli eventi incompatibili, quindi, \textbf{le probabilità si sommano}.

\begin{esempio}{L'unione e il complementare al lavoro}
\begin{itemize}
  \item Un dado onesto: la probabilità di ottenere un numero pari \emph{oppure} almeno 5 è
        \[ P(\{2,4,6\} \cup \{5,6\}) = \frac{3}{6} + \frac{2}{6} - \frac{1}{6} = \frac{4}{6} = \frac{2}{3} , \]
        dove si toglie $P(\{6\}) = 1/6$ perché il 6 è sia pari sia almeno 5. In effetti l'unione è $\{2, 4, 5, 6\}$.
  \item Due lanci di un dado: «almeno un 6» è l'unione di «6 al primo lancio» e «6 al secondo», che si verificano
        insieme solo con $(6,6)$:
        \[ P(\text{almeno un 6}) = \frac{6}{36} + \frac{6}{36} - \frac{1}{36} = \frac{11}{36} , \]
        lo stesso numero trovato contando le caselle nel capitolo precedente. Ancora più rapido è passare al
        complementare: «nessun 6» vuol dire 5 possibilità per ogni lancio, $5 \cdot 5 = 25$ coppie, e
        $P(\text{almeno un 6}) = 1 - \frac{25}{36} = \frac{11}{36}$.
\end{itemize}
\end{esempio}

\info{Il trucco del complementare}{
  Quando un evento contiene la parola «almeno», spesso conviene calcolare la probabilità del suo complementare, che
  contiene la parola «nessuno», e sottrarla da 1. Con dieci lanci di un dado, per esempio, «almeno un 6» richiederebbe
  di considerare un 6, due 6, \ldots, dieci 6; «nessun 6» invece è un solo conto: $5^{10}$ sequenze su $6^{10}$, e
  $P(\text{almeno un 6}) = 1 - (5/6)^{10} \approx 0{,}84$.
}

% ============================================================
\section{Probabilità condizionata}
% ============================================================

Spesso si ha un'informazione parziale sul risultato: «è uscito un numero pari», «il pacchetto è arrivato in
ritardo». Come cambia la probabilità degli altri eventi, sapendo questo?

Siano $A$ e $B$ due eventi che, in $n$ esperimenti, si verificano $n_A$ e $n_B$ volte ($n_B > 0$). Definiamo la
\term{frequenza di $A$ condizionata a $B$}, $f_n(A \mid B)$, come il rapporto fra il numero di volte in cui si
verificano \emph{entrambi} ($n_{A \cap B}$) e il numero di volte in cui si verifica $B$:
\[
  f_n(A \mid B) = \frac{n_{A \cap B}}{n_B} = \frac{n_{A \cap B}}{n} \cdot \frac{n}{n_B}
  = \frac{f_n(A \cap B)}{f_n(B)} .
\]
In altre parole, \textbf{restringiamo il nostro campione} (il nostro «mondo») ai soli $n_B$ esperimenti in cui si è
verificato $B$, e fra questi guardiamo in quanti si è verificato anche $A$. Passando al limite (il rapporto dei limiti
si può fare perché $P(B) > 0$):

\dfn{Probabilità condizionata}{
  Se $P(B) > 0$, la \term{probabilità di $A$ condizionata a $B$} (o «di $A$ dato $B$») è
  \[ P(A \mid B) = \frac{P(A \cap B)}{P(B)} . \]
}

\begin{figure}[H]
\centering
\venn{\fill[mspigray!30, even odd rule] \vennOmega \vennB;
      \begin{scope}\clip \vennB; \fill[evid] \vennA; \end{scope}}{si scarta tutto ciò che sta fuori da $B$}
\hspace{1.5cm}
\begin{tikzpicture}[baseline=(current bounding box.north)]
  \fill[mspiB!8] (0,0) circle (0.95);
  \begin{scope}\clip (0,0) circle (0.95); \fill[evid] (-0.8,0) circle (0.95); \end{scope}
  \draw[mspiB!85!black, line width=0.8pt] (0,0) circle (0.95);
  \node[font=\small] at (-0.42,0) {$A \cap B$};
  \node[mspiB!85!black, font=\small] at (0.85,0.85) {$B$};
  \node[font=\small, anchor=north, text width=4.2cm, align=center] at (0,-1.1)
    {sapendo che è successo $B$: il nuovo mondo è $B$, e di $A$ resta solo $A \cap B$};
\end{tikzpicture}
\caption{La probabilità condizionata: si scarta tutto ciò che sta fuori da $B$ (in grigio) e si misura quanta parte di
$B$ è occupata da $A$.}
\end{figure}

\begin{esempio}{Condizionare un dado}
Si lancia un dado onesto.
\begin{itemize}
  \item Sapendo che è uscito un numero pari, qual è la probabilità che sia il 6?
        \[ P(\{6\} \mid \text{pari}) = \frac{P(\{6\} \cap \{2,4,6\})}{P(\{2,4,6\})} = \frac{1/6}{3/6} = \frac{1}{3} . \]
        Ha senso: fra i tre pari, il 6 è uno.
  \item Sapendo che è uscito un numero non superiore a 3, qual è la probabilità che sia dispari?
        \[ P(\text{dispari} \mid \{1,2,3\}) = \frac{P(\{1,3\})}{P(\{1,2,3\})} = \frac{2/6}{3/6} = \frac{2}{3} . \]
  \item Viceversa $P(\text{pari} \mid \{6\}) = 1$: se è uscito 6, è certamente uscito un pari. Il condizionamento
        \textbf{non è simmetrico}: $P(A \mid B)$ e $P(B \mid A)$ sono, in generale, numeri diversi.
\end{itemize}
\end{esempio}

\subsection{La legge della probabilità composta}

Scambiando i ruoli di $A$ e $B$, la stessa definizione dà
\[ P(B \mid A) = \frac{P(A \cap B)}{P(A)} , \qquad P(A) > 0 . \]
Moltiplicando ciascuna delle due definizioni per il suo denominatore si ottiene una formula per l'intersezione.

\begin{formula}[legge della probabilità composta]
\[ P(A \cap B) = P(B \mid A)\, P(A) = P(A \mid B)\, P(B) \]
\end{formula}

È utilissima quando un esperimento avviene \textbf{in due tempi}: la probabilità che succedano entrambe le cose è la
probabilità della prima, moltiplicata per la probabilità della seconda \emph{sapendo} che la prima è successa.

\begin{esempio}{Due assi di fila}
Da un mazzo di 52 carte se ne estraggono due, senza reinserimento. Qual è la probabilità che siano due assi? Siano
$A_1$ = «la prima è un asso» e $A_2$ = «la seconda è un asso». Alla prima estrazione gli assi sono 4 su 52; se la
prima è stata un asso, alla seconda restano 3 assi su 51 carte:
\[ P(A_1 \cap A_2) = P(A_1)\, P(A_2 \mid A_1) = \frac{4}{52} \cdot \frac{3}{51} = \frac{12}{2\,652} = \frac{1}{221}
   \approx 0{,}0045 . \]
\begin{center}
\begin{tikzpicture}[x=1cm, y=0.62cm, every node/.style={font=\small}]
  \node[draw, circle, fill=mspilight, inner sep=2pt] (r) at (0,0) {};
  \node[draw, rounded corners=2pt, fill=mspiA!12] (a) at (3,1.5) {asso};
  \node[draw, rounded corners=2pt, fill=mspilight] (n) at (3,-1.5) {non asso};
  \draw (r) -- node[above, sloped, font=\scriptsize] {$4/52$} (a);
  \draw (r) -- node[below, sloped, font=\scriptsize] {$48/52$} (n);
  \node[draw, rounded corners=2pt, fill=evid] (aa) at (6.6,2.3) {asso};
  \node[draw, rounded corners=2pt, fill=mspilight] (an) at (6.6,0.7) {non asso};
  \draw (a.east) -- node[above, sloped, font=\scriptsize] {$3/51$} (aa.west);
  \draw (a.east) -- node[below, sloped, font=\scriptsize] {$48/51$} (an.west);
  \node[anchor=west] at (7.4,2.3) {$\dfrac{4}{52} \cdot \dfrac{3}{51} = \dfrac{1}{221}$};
  \node[font=\scriptsize\bfseries] at (3,3.4) {1ª carta};
  \node[font=\scriptsize\bfseries] at (6.6,3.4) {2ª carta};
  \node[mspigray, font=\scriptsize] at (6.0,-1.5) {(rami non disegnati)};
\end{tikzpicture}
\end{center}
Lo stesso risultato si ottiene contando: le coppie di assi sono $\binom{4}{2} = 6$ su $\binom{52}{2} = 1\,326$ coppie
possibili, e $6/1\,326 = 1/221$. Nell'albero la probabilità di un percorso è il prodotto delle probabilità dei suoi
rami: è la legge della probabilità composta disegnata.
\end{esempio}

% ============================================================
\section{La legge della probabilità totale}
% ============================================================

Un'importante conseguenza della probabilità condizionata permette di calcolare la probabilità di un evento
«spezzandolo» in casi.

\dfn{Partizione}{
  Gli eventi $E_1, E_2, \ldots, E_M$, cioè $\{E_i\}_{i=1}^{M}$, formano una \term{partizione} di $\Omega$ se
  \textbf{coprono} tutto $\Omega$ e sono \textbf{a due a due incompatibili}:
  \[ \bigcup_{i=1}^{M} E_i = \Omega , \qquad E_i \cap E_j = \emptyset \quad \text{per } i \neq j . \]
  In ogni prova si verifica uno e un solo $E_i$: sono i «casi» in cui si divide il mondo.
}

\begin{figure}[H]
\centering
\begin{tikzpicture}
  \foreach \i/\col in {0/mspiA, 1/mspiB, 2/mspiC, 3/mspiric} {
    \fill[\col!12] ({-3.6+\i*1.8},-1.4) rectangle ({-1.8+\i*1.8},1.4);
    \node[\col!85!black, font=\small] at ({-2.7+\i*1.8},1.12) {$E_{\the\numexpr\i+1\relax}$};
  }
  \begin{scope}
    \clip (0,-0.1) ellipse (3.25 and 0.85);
    \fill[evid] (-3.6,-1.4) rectangle (3.6,1.4);
  \end{scope}
  \foreach \x in {-1.8,0,1.8} \draw[mspigray, dashed] (\x,-1.4) -- (\x,1.4);
  \draw[black, line width=0.8pt] (0,-0.1) ellipse (3.25 and 0.85);
  \foreach \x/\i in {-2.5/1, -0.9/2, 0.9/3, 2.5/4}
    \node[font=\scriptsize] at (\x,-0.1) {$A \cap E_\i$};
  \node[font=\small] at (3.35,0.7) {$A$};
  \draw[mspigray, line width=0.5pt, rounded corners=2pt] (-3.6,-1.4) rectangle (3.6,1.4);
  \node[mspigray, font=\footnotesize, anchor=south east] at (3.6,-1.4) {$\Omega$};
\end{tikzpicture}
\caption{Una partizione di $\Omega$ in quattro casi. L'evento $A$ viene tagliato in quattro pezzi $A \cap E_i$, che
non si sovrappongono: la probabilità di $A$ è la somma delle probabilità dei pezzi.}
\end{figure}

\teorema[totale]{Legge della probabilità totale}{
  Se $\{E_i\}_{i=1}^{M}$ è una partizione di $\Omega$ e $A \subseteq \Omega$ è un evento qualsiasi, allora
  \[ P(A) = \sum_{i=1}^{M} P(A \cap E_i) = \sum_{i=1}^{M} P(A \mid E_i)\, P(E_i) . \]
}

\begin{dimostrazione}
\passo{Si «ritaglia» $A$ con la partizione} $A$ è contenuto in $\Omega$, quindi $A = A \cap \Omega$; e $\Omega$ è
l'unione degli $E_i$. Per la proprietà distributiva (Teorema~\ref{teo:distributiva}):
\[ P(A) = P(A \cap \Omega) = P\left( A \cap \bigcup_{i=1}^{M} E_i \right) = P\left( \bigcup_{i=1}^{M} (A \cap E_i)
   \right) . \]
\passo{I pezzi non si sovrappongono} Siccome $E_i \cap E_j = \emptyset$ per $i \neq j$, anche i pezzi di $A$ sono
incompatibili: $(A \cap E_i) \cap (A \cap E_j) \subseteq E_i \cap E_j = \emptyset$.
\passo{Si sommano le probabilità} Per eventi incompatibili la probabilità dell'unione è la somma delle probabilità
(proprietà b) con $P(A \cap B) = 0$, applicata un pezzo alla volta):
\[ P(A) = \sum_{i=1}^{M} P(A \cap E_i) . \]
\passo{Si usa la probabilità composta} Ogni pezzo si riscrive come $P(A \cap E_i) = P(A \mid E_i)\, P(E_i)$, e si
ottiene la seconda forma.
\end{dimostrazione}

\begin{esempio}{Un canale binario}
Una sorgente trasmette il bit 0 con probabilità $0{,}6$ e il bit 1 con probabilità $0{,}4$. Il canale è rumoroso:
qualunque bit venga trasmesso, con probabilità $0{,}1$ arriva invertito. Qual è la probabilità di ricevere un 1?
\begin{center}
\begin{tikzpicture}[bit/.style={draw, circle, minimum size=0.75cm, font=\small}]
  \node[bit, fill=mspiA!12] (t0) at (0,1) {0}; \node[bit, fill=mspiA!12] (t1) at (0,-1) {1};
  \node[bit, fill=mspiB!12] (r0) at (5,1) {0}; \node[bit, fill=mspiB!12] (r1) at (5,-1) {1};
  \node[font=\small, anchor=east] at (-0.5,1) {$P(0\ \text{tr.}) = 0{,}6$};
  \node[font=\small, anchor=east] at (-0.5,-1) {$P(1\ \text{tr.}) = 0{,}4$};
  \draw[-{Stealth}, line width=0.8pt] (t0) -- node[above, font=\scriptsize] {$0{,}9$} (r0);
  \draw[-{Stealth}, line width=0.8pt] (t1) -- node[below, font=\scriptsize] {$0{,}9$} (r1);
  \draw[-{Stealth}, line width=0.6pt, mspiB] (t0) -- node[pos=0.25, above, font=\scriptsize] {$0{,}1$} (r1);
  \draw[-{Stealth}, line width=0.6pt, mspiB] (t1) -- node[pos=0.25, below, font=\scriptsize] {$0{,}1$} (r0);
  \node[font=\scriptsize\bfseries] at (0,1.9) {trasmesso};
  \node[font=\scriptsize\bfseries] at (5,1.9) {ricevuto};
\end{tikzpicture}
\end{center}
I due casi «è stato trasmesso 0» ed «è stato trasmesso 1» formano una partizione. Per la legge della probabilità
totale:
\[
  P(1\ \text{ric.}) = P(1\ \text{ric.} \mid 0\ \text{tr.})\, P(0\ \text{tr.}) + P(1\ \text{ric.} \mid 1\ \text{tr.})\,
  P(1\ \text{tr.}) = 0{,}1 \cdot 0{,}6 + 0{,}9 \cdot 0{,}4 = 0{,}06 + 0{,}36 = 0{,}42 .
\]
Un 1 si può ricevere in due modi: perché è stato trasmesso 1 ed è arrivato giusto, oppure perché è stato trasmesso 0
ed è stato invertito. La legge somma i due modi, ciascuno pesato con la probabilità del suo caso. Nel prossimo
capitolo ci chiederemo il contrario: \emph{ricevuto} un 1, qual è la probabilità che sia stato \emph{trasmesso} un 1?
\end{esempio}

% ============================================================
\section{Eventi indipendenti}
% ============================================================

Due eventi $A$ e $B$ si dicono \textbf{indipendenti} quando il verificarsi di uno non ha alcuna influenza sul
verificarsi o meno dell'altro. Con le frequenze: sapere che si è verificato $B$ non cambia la frequenza di $A$, e
viceversa,
\[ f_n(A \mid B) \simeq f_n(A) \qquad\text{e}\qquad f_n(B \mid A) \simeq f_n(B) . \]
Usando la legge della probabilità composta, in questo caso
\[ f_n(A \cap B) = f_n(A)\, f_n(B \mid A) = f_n(B)\, f_n(A \mid B) \simeq f_n(A)\, f_n(B) , \]
e passando al limite si arriva alla definizione.

\dfn{Eventi indipendenti}{
  Due eventi $A$ e $B$ sono \term{indipendenti} se (e solo se)
  \[ P(A \cap B) = P(A)\, P(B) . \]
}

\teorema[indip]{Indipendenza e probabilità condizionata}{
  Se $P(B) > 0$, gli eventi $A$ e $B$ sono indipendenti se e solo se $P(A \mid B) = P(A)$. Allo stesso modo, se
  $P(A) > 0$, sono indipendenti se e solo se $P(B \mid A) = P(B)$.
}

\begin{dimostrazione}
\passo{Da indipendenti a condizionata} Se $P(A \cap B) = P(A)\,P(B)$, allora
\[ P(A \mid B) = \frac{P(A \cap B)}{P(B)} = \frac{P(A)\,P(B)}{P(B)} = P(A) . \]
\passo{Da condizionata a indipendenti} Se $P(A \mid B) = P(A)$, per la probabilità composta
$P(A \cap B) = P(A \mid B)\,P(B) = P(A)\,P(B)$.
\passo{Scambiando i ruoli} Lo stesso ragionamento, con $A$ e $B$ scambiati, vale per $P(B \mid A)$.
\end{dimostrazione}

La forma con il prodotto è preferibile come definizione perché è simmetrica e funziona anche quando una delle due
probabilità è nulla (la probabilità condizionata a un evento di probabilità zero non è definita).

\begin{esempio}{Due lanci di un dado}
I due lanci di un dado onesto sono indipendenti: il primo non può influenzare il secondo. Per ogni coppia $(i, j)$:
\[ P(\{(i, j)\}) = P(\text{primo} = i)\, P(\text{secondo} = j) = \frac{1}{6} \cdot \frac{1}{6} = \frac{1}{36} , \]
che è proprio la probabilità di ciascuna delle 36 coppie equiprobabili del capitolo precedente. Allo stesso modo
\[ P(\text{primo pari e secondo} = 6) = \frac{1}{2} \cdot \frac{1}{6} = \frac{1}{12} , \]
e in effetti le coppie favorevoli sono $(2,6)$, $(4,6)$, $(6,6)$, cioè 3 su 36.
\end{esempio}

\begin{esempio}{Indipendenti anche senza «separazione fisica»}
In un \textbf{solo} lancio di un dado onesto, siano $A = \{2, 4, 6\}$ (pari) e $B = \{1, 2\}$ (non più di 2). Allora
$A \cap B = \{2\}$ e
\[ P(A \cap B) = \frac{1}{6} = \frac{1}{2} \cdot \frac{1}{3} = P(A)\,P(B) : \]
i due eventi sono indipendenti. Sapere che è uscito 1 o 2 non cambia la probabilità che sia uscito un pari: è $1/2$
prima ed è $1/2$ dopo (uno dei due casi rimasti è pari). L'indipendenza è una proprietà dei \emph{numeri}, non solo
della «fisica» dell'esperimento.
\end{esempio}

\warning{Indipendenti non vuol dire incompatibili}{
  È un errore molto comune confondere le due parole, che dicono quasi il contrario. Se $A$ e $B$ sono
  \textbf{incompatibili} ($A \cap B = \emptyset$) e hanno entrambi probabilità positiva, allora
  $P(A \cap B) = 0 \neq P(A)\,P(B)$: sono \textbf{dipendenti}, e in modo fortissimo, perché se si verifica uno
  sicuramente l'altro non si verifica. Testa e croce nello stesso lancio sono incompatibili; testa al primo lancio e
  testa al secondo sono indipendenti.
}

\gbox{In sintesi}{mspiA!80!black}{
\begin{itemize}
  \item La \textbf{probabilità} di un evento è il limite della sua frequenza relativa: $P(A) = \lim_{n\to\infty}
        n_A/n$. Se gli eventi elementari sono equiprobabili, $P(A) = \card{A}/\card{\Omega}$ (casi favorevoli su casi
        possibili), purché i casi contati siano davvero equiprobabili.
  \item Se non lo sono, eventi con lo stesso numero di elementi possono avere probabilità diverse (dado truccato).
  \item Regole: $P(\compl{A}) = 1 - P(A)$; $P(A \cup B) = P(A) + P(B) - P(A \cap B)$;
        $P(A \setminus B) = P(A) - P(A \cap B)$; $P(\Omega) = 1$, $P(\emptyset) = 0$; per eventi incompatibili le
        probabilità si sommano.
  \item Probabilità condizionata: $P(A \mid B) = P(A \cap B)/P(B)$, cioè si restringe il mondo a $B$. Probabilità
        composta: $P(A \cap B) = P(B \mid A)\,P(A) = P(A \mid B)\,P(B)$.
  \item Probabilità totale: se gli $E_i$ formano una partizione, $P(A) = \sum_i P(A \mid E_i)\,P(E_i)$.
  \item $A$ e $B$ sono indipendenti se $P(A \cap B) = P(A)\,P(B)$, cioè se sapere che uno si è verificato non cambia
        la probabilità dell'altro. Indipendenti non vuol dire incompatibili.
\end{itemize}
}
